\section{Service Provisioning}
\label{chap:proposal:provisioning}

To implement and validate network applications, third parties require a network environment that reproduces the conditions under which each application is supposed to operate. The role of a testbed is not only to provide such a substrate, i.e. the actual physical network and compute infrastructure, but also to provide the necessary capabilities to configure it for each experiment, according to the requirements of the application and target vertical industry. This requires partitioning the shared infrastructure into isolated logical networks through network slicing, exposed as third-party services that let verticals configure slices tailored to their applications. Moreover, the complete lifecycle of the provisioned slices and their UE should follow widely adopted, accessible data models.

Composing complete experiments, however, requires access to other functionalities. The exposure APIs, monitoring services, and the vertical application itself must all be configured, packaged and deployed for each experiment. These capabilities should be offered as services through a single management and automation platform, traditionally an OSS. Through it, verticals can design, order, and manage services as needed, without manually interacting with the underlying infrastructure components which realize their orchestration. Centralizing the provisioning of experiments in a single platform would also ensure that each experiment is instantiated consistently, so that its results can be reproduced and compared across executions.

With this in mind, vertical clients should be able to access this platform both through a graphical interface and programmatically, via standardized APIs that support automation and integration with external tools. TMF ODA addresses this by defining TMF Open APIs to simplify service design, orchestration, and management, abstracting proprietary management interfaces so network capabilities are offered through standardized, interoperable ones.

The remainder of this section details how such a platform may support the complete lifecycle of an experiment by translating vertical requirements into services that can be provisioned and tested on the testbed.


\subsection{Service Design}
\label{chap:proposal:provisioning:design}

\subsection{Service Orchestration \& Lifecycle Management}
\label{chap:proposal:provisioning:orchestration-lifecycle}

\subsection{Service Validation}
\label{chap:proposal:provisioning:validation}

The orchestration of this service chain is merely the first step in the experimentation process, establishing the underlying environment for the evaluation of vertical applications. This evaluation must verify that the application functions as intended and effectively utilizes the exposed \acs{5G} capabilities, just deployed as services. In principle, this could be achieved by the experimenters themselves, who could programmatically extract monitoring metrics and application logs from the services offered by the testbed (e.g. Prometheus, Grafana or even \acs{NEF} and CAMARA \acsp{API}) and its orchestration platforms (e.g. Kubernetes, \acs{NFVO}). However, this approach places an unnecessary burden on vertical experimenters, who would have to develop their own collection and validation procedures for tasks that are largely common across experiments and could readily be automated by the testbed. Testing should therefore be configured and offered via the \acs{OSS} as part of the experiment, rather than left for vertical experimenters to conduct ad hoc once the services are active.

To this end, the \acs{OSS} should integrate with a \acs{CI} platform that automates testing procedures and expose standardized \acsp{API} through which verticals can configure and execute them. As discussed in Section~\ref{sec:pipelines}, such solutions have already been adopted in the literature to orchestrate and validate network applications over \acs{TMF}-aligned platforms~\cite{AutomatedTestingOrchestrationFramework, TMFOpenDigitalArchitecture5GServiceOrchestration}. Test cases should be declared alongside the experiments' services and configured through these \acsp{API}, with the \acs{OSS} triggering their execution as they become active. Validation should also proceed sequentially within a multi-stage pipeline: each service gets provisioned and validated before the subsequent one is instantiated, so that failures can be detected and attributed to the stage where they occur, without compromising the rest of the experiment. For instance, if a requested network slice fails to be instantiated, it should fail the test and optionally trigger a teardown of the entire service chain before the network applications that depend on it are deployed. This applies both to the services provided by the testbed, and to the network application, whose behavior is assessed against the data collected during the tests execution. At each stage, the \acs{CI} platform must collect the relevant network monitoring metrics, application and orchestration logs and exposure \acs{API} responses, and compare them against the expected values, such as the attributes of the ordered \acs{NEST} or the application \acsp{KPI}, to produce a verdict for each test. These results should then be aggregated into reports associated with the corresponding service instances and made available to verticals through the \acs{OSS}.

In this way, verticals need only to specify what must be validated in their experiments, while the testbed remains responsible for how these tests are actually implemented, orchestrated, and executed across the underlying infrastructure.
